% Generated draft: requires mathematical and source review.
\documentclass[11pt]{article}
\usepackage[margin=25mm]{geometry}
\usepackage{fontspec}
\setmainfont{Latin Modern Roman}
\usepackage{amsmath,amssymb,mathtools,unicode-math}
\setmathfont{Latin Modern Math}
\usepackage{enumitem,needspace,xurl,hyperref}
\setlength{\parindent}{0pt}
\setlength{\parskip}{5pt}
\setlength{\emergencystretch}{3em}
\newcommand{\N}{\mathbb{N}}
\newcommand{\Z}{\mathbb{Z}}
\newcommand{\Q}{\mathbb{Q}}
\newcommand{\R}{\mathbb{R}}
\newcommand{\C}{\mathbb{C}}
\newcommand{\F}{\mathbb{F}}
\newcommand{\A}{\mathbb{A}}
\newcommand{\PP}{\mathbb{P}}
\newcommand{\E}{\mathbb{E}}
\newcommand{\eps}{\varepsilon}
\newcommand{\dd}{\,\mathrm{d}}
\begin{document}
\section*{UnsolvedMath ID: 800157}
\textbf{Source problem code:} catalogue entry
\subsection{Definitions and mathematical statement}
Let \(k\in\{1,2,3,\ldots\}\) and let \(T>0\). The supplied excerpt does not define \(I_k(T)\). Under the standard interpretation suggested by its title, let \(\zeta(s)\) denote the meromorphic continuation of \(\sum_{n=1}^{\infty}n^{-s}\), initially defined for \(\operatorname{Re}s>1\), and set
\[
I_k(T)=\int_0^T\left|\zeta\!\left(\tfrac12+it\right)\right|^{2k}\,\dd t.
\]
This interpretation, including the lower integration endpoint, requires confirmation from the full source. For each \(k\), let \(P_{k^2}(x)\) be the particular polynomial of degree \(k^2\) intended in the conjecture attributed by the excerpt to Conrey et al., with reference label [CFKRS]. Its coefficients are not supplied; the polynomial must be fixed as intended by that conjecture, rather than chosen freely. Define
\[
E_k(T)=I_k(T)-T P_{k^2}(\log T).
\]
The excerpt attributes to Conrey et al. the conjectural estimate
\[
E_k(T)=O_{k,\varepsilon}\!\left(T^{1/2+\varepsilon}\right)
\qquad(T\to\infty)
\]
for each fixed \(k\) and every \(\varepsilon>0\), and reports subsequent debate about the correct size of this error term.

\textbf{Problem.} Subject to confirmation of the moment normalization and the intended polynomials, formulate a conjecture for the optimal exponents \(\theta_k\), especially for \(k\geq3\), such that, for every fixed positive integer \(k\) and every \(\varepsilon>0\), there exist constants \(C_{k,\varepsilon}>0\) and \(T_{0,k,\varepsilon}>0\) satisfying
\[
|E_k(T)|\leq C_{k,\varepsilon}T^{\theta_k+\varepsilon}
\qquad\text{for all }T\geq T_{0,k,\varepsilon}.
\]
A precise interpretation of optimality is
\[
\theta_k=\inf\left\{\alpha\in\R:
\text{for every }\varepsilon>0,\quad
E_k(T)=O_{k,\alpha,\varepsilon}\!\left(T^{\alpha+\varepsilon}\right)
\text{ as }T\to\infty\right\}.
\]
This interpretation presupposes the intended fixed polynomial and seeks a finite real exponent; it makes explicit the source's otherwise informal use of optimality.

The source reports that \(\theta_1=\tfrac14\) and \(\theta_2=\tfrac12\) are widely believed to be optimal, mentions omega results in both cases without stating them, and specifically attributes these proposed exponents to N. Ng. It reports that both Ivi\'c and Motohashi suggested \(\theta_3=\tfrac34\). These are supplied conjectural assessments, not assertions established here. The problem also asks whether the Multiple Dirichlet Series method can yield an appropriate conjecture and mentions work in progress by Baluyot and Cech studying this question. That work-in-progress description is retained as historical source context, without a claim about its present status. The source contribution is attributed to Nathan Ng, Lethbridge.
\subsection{Short English statement}
Conjecture the sharp growth exponent of the error in each even moment of the Riemann zeta function, particularly for \(k\geq3\), and investigate whether Multiple Dirichlet Series can predict it. The supplied suggestions are \(\theta_1=1/4\), \(\theta_2=1/2\), and \(\theta_3=3/4\). The moment normalization and the specific main-term polynomials require confirmation.
\subsection{Sources}
\begin{itemize}
\item Supplied local UnsolvedMath registry record: permanent ID \(800157\), source problem code \texttt{catalogue entry}, title ``Moments of L-functions --- Moments of the Riemann zeta function --- Problem 6.'' The registry attributes the problem to Hamieh, Kadiri, and Ng, with proposed year 2024, and labels it open. This is a supplied status label, not an independently verified mathematical assessment.
\item Hamieh, Kadiri, and Ng, \emph{Moments of L-functions Problem List} (2024), Problem 6. Supplied source URL: \url{https://arxiv.org/abs/2405.17800}. The supplied provenance records access on 2026-10-16 and extraction by \texttt{arxiv\_latex}; neither access nor the source contents have been independently verified here. The contribution is signed Nathan Ng, Lethbridge.
\item The excerpt cites Conrey et al. using the label [CFKRS], but supplies no full bibliographic entry. It also names Ivi\'c and Motohashi in connection with the proposed third-moment-error exponent, and Baluyot and Cech in connection with work in progress on Multiple Dirichlet Series; no separate references or links for these statements are supplied.
\end{itemize}
\end{document}
